\documentclass[11pt]{article}

\usepackage{amsmath,amssymb,amsthm,mathtools}
\usepackage[margin=1in]{geometry}
\usepackage{hyperref,microtype}

\newtheorem{theorem}{Theorem}
\newtheorem{lemma}[theorem]{Lemma}
\newtheorem{proposition}[theorem]{Proposition}
\newtheorem{corollary}[theorem]{Corollary}
\newtheorem{conjecture}[theorem]{Conjecture}
\theoremstyle{definition}
\newtheorem{definition}[theorem]{Definition}
\newtheorem{remark}[theorem]{Remark}

\newcommand{\eps}{\varepsilon}
\newcommand{\E}{\mathbb{E}}
\newcommand{\Var}{\mathrm{Var}}

\title{On the HHKP $\tfrac12$-conjecture\\
\large (an attempt, not a proof)}
\date{}

\begin{document}
\maketitle

\begin{abstract}
Hefty--Horn--King--Pfender and Campos--Jenssen--Michelen--Sahasrabuddhe conjecture
\[
  R(3,k)\;=\;\Bigl(\tfrac12+o(1)\Bigr)\frac{k^{2}}{\log k}.
\]
The lower bound is their theorem. The missing half is an upper bound, i.e.\ a
$\sqrt{2}$-improvement of Shearer's independence-number estimate for every
triangle-free graph. We prove the standard reduction (doubling Shearer implies
the conjecture), prove that every argument which only lower-bounds the
\emph{average} independent-set size is blocked by the Davies--Jenssen--Perkins--Roberts
occupancy theorem, and record a two-sample attempt that dies for the same reason
local algorithms fail on $G(n,d/n)$. We do \emph{not} prove the conjecture.
Improving Shearer's constant by any fixed factor remains open.
\end{abstract}

\section{The statement}

Write $\log$ for the natural logarithm.
Shearer~\cite{Shearer1983} proved that every $n$-vertex triangle-free graph of
average degree $d$ satisfies
\begin{equation}\label{eq:shearer}
  \alpha(G)\;\ge\; f(d)\,n,
  \qquad
  f(d)\;=\;\frac{d\log d-d+1}{(d-1)^{2}}\;=\;\bigl(1+o_d(1)\bigr)\frac{\log d}{d}.
\end{equation}
Optimising against the trivial bound $\alpha(G)\ge\Delta(G)$ yields
\begin{equation}\label{eq:shearer-R}
  R(3,k)\;\le\;\bigl(1+o(1)\bigr)\frac{k^{2}}{\log k}.
\end{equation}
HHKP~\cite{HHKP2025} proved $R(3,k)\ge\bigl(\tfrac12+o(1)\bigr)k^{2}/\log k$,
and with CJMS~\cite{CJMS2025} conjectured that this is the truth.

\begin{conjecture}[HHKP / CJMS]\label{conj:half}
\[
  R(3,k)\;=\;\Bigl(\tfrac12+o(1)\Bigr)\frac{k^{2}}{\log k}.
\]
\end{conjecture}

Equivalently: every $n$-vertex triangle-free graph has
$\alpha(G)\ge\bigl(1-o(1)\bigr)\sqrt{n\log n}$.
Shearer gives only the constant $1/\sqrt{2}$.

\section{What the optimisation actually says}

\begin{lemma}[Neighbourhood versus Shearer]\label{lem:opt}
Let $G$ be an $n$-vertex triangle-free graph of average degree $d$ and maximum
degree $\Delta$. Then
\[
  \alpha(G)
  \;\ge\;
  \max\Bigl(\Delta,\; f(d)\,n\Bigr).
\]
In particular
\[
  \alpha(G)
  \;\ge\;
  \bigl(1-o(1)\bigr)\,
  \max\Bigl(\Delta,\; n\frac{\log d}{d}\Bigr).
\]
\end{lemma}

\begin{proof}
The neighbourhood of a maximum-degree vertex is independent, because $G$ is
triangle-free. The second bound is~\eqref{eq:shearer}.
\end{proof}

\begin{proposition}[Shearer's Ramsey constant]\label{prop:opt}
Write $d=\theta\sqrt{n\log n}$ with $\theta>0$ bounded away from $0$ and
$\infty$, and assume $d=\Theta(\Delta)$. Then Lemma~\ref{lem:opt} gives
\[
  \frac{\alpha(G)}{\sqrt{n\log n}}
  \;\ge\;
  \bigl(1-o(1)\bigr)\,
  \max\Bigl(\theta,\;\tfrac{1}{2\theta}\Bigr).
\]
The right-hand side is minimised at $\theta=1/\sqrt{2}$, with value $1/\sqrt{2}$.
This is~\eqref{eq:shearer-R}.
\end{proposition}

\begin{proof}
$\log d=\tfrac12\log n+\log\theta+\tfrac12\log\log n$, so
\[
  n\frac{\log d}{d}
  \;=\;
  \frac{\sqrt{n\log n}}{\theta}
  \cdot
  \frac{\tfrac12\log n+O(\log\log n)}{\log n}
  \;=\;
  \bigl(\tfrac{1}{2\theta}+o(1)\bigr)\sqrt{n\log n}.
\]
Minimise $\max(\theta,1/(2\theta))$.
\end{proof}

The HHKP graphs realise $\theta=1+o(1)$ and $\alpha=\bigl(1+o(1)\bigr)\sqrt{n\log n}$:
both a neighbourhood and a random-like independent set meet the independence
number. One cannot beat Shearer on those graphs by ``take a neighbourhood''.
The only remaining gain is a better function $f$.

\section{Doubling Shearer implies the conjecture}

\begin{conjecture}[Shearer $\times 2$]\label{conj:double}
Every $n$-vertex triangle-free graph of maximum degree $d$ satisfies
\[
  \alpha(G)\;\ge\;\bigl(2-o_d(1)\bigr)\,n\frac{\log d}{d}.
\]
\end{conjecture}

Random $d$-regular graphs (which are triangle-free with positive probability for
fixed $d$) have $\alpha=\bigl(2-o_d(1)\bigr)n\log d/d$, so the constant $2$ is
best possible if it holds. Davies--Jenssen--Perkins--Roberts~\cite{DJPR2018}
proved that the \emph{average} independent-set size is already
$\bigl(1+o_d(1)\bigr)n\log d/d$, and that the constant $1$ is tight for the
average. Conjecture~\ref{conj:double} is therefore a statement about the gap
between the maximum and the average, not about the average itself.

\begin{theorem}[Reduction]\label{thm:reduction}
Conjecture~\ref{conj:double} implies Conjecture~\ref{conj:half}.
\end{theorem}

\begin{proof}
Let $G$ be $n$-vertex triangle-free of maximum degree $d$.
Conjecture~\ref{conj:double} and $\alpha\ge\Delta\ge$ a typical degree give,
in the regular case $d=\theta\sqrt{n\log n}$,
\[
  \frac{\alpha(G)}{\sqrt{n\log n}}
  \;\ge\;
  \bigl(1-o(1)\bigr)\,
  \max\Bigl(\theta,\;\tfrac{1}{\theta}\Bigr)
  \;\ge\;
  1.
\]
Thus $\alpha(G)\ge\bigl(1-o(1)\bigr)\sqrt{n\log n}$.
If $G$ has no independent set of size $k$, then $n\le\bigl(\tfrac12+o(1)\bigr)k^{2}/\log k$.
The irregular case is the same after deleting vertices of degree
$\le d/(2\log d)$ as in~\cite[Lemma~1]{DJPR2018}: the deleted stars contribute
$o\bigl(n\log d/d\bigr)$ and the remaining graph is almost regular of degree
$\Theta(d)$.
\end{proof}

The same computation with a factor $c\in[1,2]$ in place of $2$ produces
$R(3,k)\le\bigl(1/c+o(1)\bigr)k^{2}/\log k$.
In particular, any fixed improvement of Shearer's leading constant improves the
Ramsey constant. CJMS record that even this is open:
``it is a notoriously difficult open problem to improve Shearer's bound by any
constant factor.''

\section{The occupancy barrier}

Let $I$ be a random independent set of $G$ in the hard-core model at fugacity
$\lambda>0$, and write $\alpha_G(\lambda)=\E|I|$.
Davies--Jenssen--Perkins--Roberts proved
\begin{equation}\label{eq:occup}
  \frac{\alpha_G(\lambda)}{n}
  \;\ge\;
  \frac{\lambda}{1+\lambda}
  \cdot
  \frac{W\bigl(d\log(1+\lambda)\bigr)}{d\log(1+\lambda)}
\end{equation}
for every triangle-free graph of maximum degree $d$, where $W$ is the Lambert
function. For every $\lambda=d^{o(1)}$ this is
$\bigl(1-o_d(1)\bigr)\log d/d$, and the constant $1$ is tight.

\begin{theorem}[Occupancy cannot prove Conjecture~\ref{conj:half}]\label{thm:barrier}
Any lower bound on $\alpha(G)$ that uses only $\alpha(G)\ge\alpha_G(\lambda)$ for
some $\lambda=d^{o(1)}$ yields at most Shearer's constant $1+o(1)$ in~\eqref{eq:shearer},
hence at most~\eqref{eq:shearer-R}.
\end{theorem}

\begin{proof}
For $\lambda=d^{o(1)}$,~\eqref{eq:occup} and its matching random-regular upper
bound give $\alpha_G(\lambda)=\bigl(1+o_d(1)\bigr)n\log d/d$.
The inequality $\alpha(G)\ge\alpha_G(\lambda)$ therefore cannot beat~\eqref{eq:shearer}.
\end{proof}

This includes: the uniform average over all independent sets ($\lambda=1$);
Shearer's randomised greedy (which samples a typical independent set);
every hard-core argument at bounded or slowly growing fugacity;
and every local algorithm whose output distribution is contiguous with a
bounded-fugacity hard-core sample.
The algorithmic factor-of-two gap on $G(n,d/n)$ --- no local algorithm finds
an independent set of size $(1+\eps)n\log d/d$ --- is the same obstruction
with the triangle-free hypothesis removed~\cite{GS14,RV17}.

To see the maximum one must take $\lambda\to\infty$.
The identity
\begin{equation}\label{eq:var}
  \alpha(G)
  \;=\;
  \alpha_G(1)
  \;+\;
  \int_{1}^{\infty}
  \frac{\Var_{\lambda}(|I|)}{\lambda}\,d\lambda
\end{equation}
converts Conjecture~\ref{conj:double} into a tail/variance problem:
the integral must contribute another $\bigl(1-o_d(1)\bigr)n\log d/d$.
At $\lambda=1$ the measure is concentrated on typical sets of size
$\sim n\log d/d$. The independent sets of size $\sim 2n\log d/d$ are
invisible until $\lambda$ is large enough to reweight the tail.

\begin{remark}
Davies--Jenssen--Perkins--Roberts conjectured $\alpha(G)/\alpha_G(1)\ge 2-o_d(1)$
for triangle-free graphs of minimum degree $d$, and noted that this implies
Conjecture~\ref{conj:half}. Their weaker conjecture $\alpha/\alpha_G(1)\ge 4/3$
would give only $R(3,k)\le\bigl(3/4+o(1)\bigr)k^{2}/\log k$.
The Moon--Moser bound $\alpha/\alpha_G(1)\ge 1+\alpha/n$ is $1+o(1)$ on
Ramsey graphs, so it does not help.
\end{remark}

\section{A two-sample attempt, and why it fails}

The natural try, suggested by the two-bite construction, is to draw two
independent hard-core samples $I_1,I_2$ at $\lambda=1$ and clean the bipartite
graph of crossing edges.

\begin{lemma}[Cleaning two samples]\label{lem:twosample}
Let $I_1,I_2$ be independent sets. The crossing graph $H$ with parts
$(I_1\setminus I_2,\,I_2\setminus I_1)$ is bipartite, hence triangle-free, and
\[
  \alpha(G)
  \;\ge\;
  |I_1\cap I_2|
  \;+\;
  \alpha(H).
\]
\end{lemma}

\begin{proof}
$I_1\cap I_2$ is independent of both sides of $H$ (any edge from the intersection
into $I_j$ would be an edge of $G$ inside $I_j$). A maximum independent set of
$H$ may therefore be added to the intersection.
\end{proof}

If $I_1,I_2$ were independent uniform samples of size $s=n\log d/d$ in a
$d$-regular triangle-free graph, a first-moment guess is
$|I_1\cap I_2|\sim s^{2}/n=n(\log d)^{2}/d^{2}=o(s)$ and
$e(H)\sim s^{2}\cdot(d/n)=s\log d$.
The bipartite graph $H$ then has average degree $\sim\log d$, so Caro--Wei
gives only $\alpha(H)\ge 2s/(\log d+1)=o(s)$.
Shearer on $H$ gives
\[
  \alpha(H)
  \;\ge\;
  \bigl(1-o(1)\bigr)\,
  2s\cdot\frac{\log\log d}{\log d}
  \;=\;
  o(s).
\]
One does not recover a second copy of $s$. The crossing graph is too dense:
two typical independent sets of size $n\log d/d$ behave like a random bipartite
graph of degree $\log d$, whose independence number is $o(s)$.

Taking $I_2$ to be a neighbourhood does not help on the HHKP examples, where
neighbourhoods already have size $\alpha(G)$.
Taking $\lambda$ large enough that $\E|I|\ge(1+\eps)n\log d/d$ is exactly
Conjecture~\ref{conj:double} again.

\section{What a proof must do}

A proof of Conjecture~\ref{conj:half} must produce, in every triangle-free
graph, an independent set that is atypically large for the hard-core measure
at bounded fugacity. It cannot be:
\begin{itemize}
  \item a neighbourhood (already used in Lemma~\ref{lem:opt});
  \item a bounded-fugacity hard-core sample (Theorem~\ref{thm:barrier});
  \item a local algorithm of the $G(n,d/n)$ type;
  \item a second typical sample cleaned against the first (Lemma~\ref{lem:twosample}).
\end{itemize}
The missing input is a uniform tail bound: either
$\Var_{\lambda}(|I|)$ large on a long $\lambda$-interval, or a
non-local construction of an independent set of size
$(2-\eps)n\log d/d$. We do not have either.

\section{What is not claimed}

We do not prove Conjecture~\ref{conj:half}, Conjecture~\ref{conj:double}, or
any improvement of Shearer's leading constant.
We do not claim $R(3,k)\le\bigl(c+o(1)\bigr)k^{2}/\log k$ for any $c<1$.
Family~A and the counting derandomization of HHKP are lower-bound constructions;
they cannot prove an upper bound.

The theorems in this note are the reduction (Theorem~\ref{thm:reduction}),
the occupancy barrier (Theorem~\ref{thm:barrier}), and the failure of the
two-sample cleaning (Lemma~\ref{lem:twosample} plus the degree computation).

\begin{thebibliography}{9}
\bibitem{CJMS2025}
M.~Campos, M.~Jenssen, M.~Michelen and J.~Sahasrabuddhe,
\emph{A new lower bound for the Ramsey numbers $R(3,k)$},
arXiv:2505.13371, 2025.

\bibitem{DJPR2018}
E.~Davies, M.~Jenssen, W.~Perkins and B.~Roberts,
\emph{On the average size of independent sets in triangle-free graphs},
Proc.\ Amer.\ Math.\ Soc.\ 146 (2018), 111--124.

\bibitem{GS14}
D.~Gamarnik and M.~Sudan,
\emph{Limits of local algorithms over sparse random graphs},
Ann.\ Probab.\ 45 (2017), no.~4, 2353--2376.
\newblock arXiv:1304.1831.

\bibitem{HHKP2025}
Z.~Hefty, P.~Horn, D.~King and F.~Pfender,
\emph{Improving $R(3,k)$ in just two bites},
arXiv:2510.19718, 2025.

\bibitem{RV17}
M.~Rahman and B.~Vir\'ag,
\emph{Local algorithms for independent sets are half-optimal},
Ann.\ Probab.\ 45 (2017), no.~3, 1543--1577.

\bibitem{Shearer1983}
J.~B.~Shearer,
\emph{A note on the independence number of triangle-free graphs},
Discrete Math.\ 46 (1983), 83--87.
\end{thebibliography}

\end{document}
